% Generated draft: requires mathematical and source review.
\documentclass[11pt]{article}
\usepackage[margin=25mm]{geometry}
\usepackage{fontspec}
\setmainfont{Latin Modern Roman}
\usepackage{amsmath,amssymb,mathtools,unicode-math}
\setmathfont{Latin Modern Math}
\usepackage{enumitem,needspace,xurl,hyperref}
\setlength{\parindent}{0pt}
\setlength{\parskip}{5pt}
\setlength{\emergencystretch}{3em}
\newcommand{\N}{\mathbb{N}}
\newcommand{\Z}{\mathbb{Z}}
\newcommand{\Q}{\mathbb{Q}}
\newcommand{\R}{\mathbb{R}}
\newcommand{\C}{\mathbb{C}}
\newcommand{\F}{\mathbb{F}}
\newcommand{\A}{\mathbb{A}}
\newcommand{\PP}{\mathbb{P}}
\newcommand{\E}{\mathbb{E}}
\newcommand{\eps}{\varepsilon}
\newcommand{\dd}{\,\mathrm{d}}
\begin{document}
\section*{UnsolvedMath ID: 30004782}
\textbf{Source problem code:} catalogue entry
\subsection{Definitions and mathematical statement}
Let \(k\) range over number fields, let \(\mathbb A_k\) denote its ring of adeles, and let \(G\) be a split reductive algebraic group over \(k\). Let \(\widehat G\) denote the complex dual group of \(G\). Let \(\sigma\) be a cuspidal automorphic representation of \(G(\mathbb A_k)\), and let \(\rho\colon\widehat G\to\operatorname{GL}(V)\) be a finite-dimensional dual-group representation. Write \(\widetilde\sigma\) for the contragredient representation of \(\sigma\).

At an unramified finite place \(v\), let \(q_v\) be the residue-field cardinality and let \(c_v(\sigma)\) denote the Satake conjugacy class, once a normalization has been fixed. The corresponding unramified local factor is
\[
L_v(s,\sigma,\rho)=\det\!\left(I_V-q_v^{-s}\rho(c_v(\sigma))\right)^{-1}.
\]
For a finite set \(S\) containing the archimedean places and all ramified places, these factors specify the partial Euler product
\[
L^S(s,\sigma,\rho)=\prod_{v\notin S}L_v(s,\sigma,\rho)
\]
where it converges. The supplied record does not specify the normalization of these factors, the factors at the omitted places, or the epsilon factor.

\textbf{Requested assertion, subject to the unresolved conventions below.} For every tuple \((k,G,\sigma,\rho)\) as above, the associated automorphic \(L\)-function \(L(s,\sigma,\rho)\), initially defined in its domain of convergence, admits meromorphic continuation to all \(s\in\C\) and satisfies the identity of meromorphic functions
\[
L(s,\sigma,\rho)=\epsilon(s,\sigma,\rho)\,L(1-s,\widetilde\sigma,\rho).
\]
The same representation \(\rho\) on both sides is retained exactly as supplied.

\textbf{Unresolved qualifications.} The record does not say whether \(L\) is completed or has omitted local factors; the displayed functional equation therefore requires a compatible specification of all local factors and normalizations. Its epsilon factor must be specified independently as part of the intended automorphic construction, rather than defined as the quotient of the two displayed \(L\)-functions. The record also does not explicitly specify connectedness in the term ``split reductive group,'' the representation category intended for \(\rho\), or any conventions concerning central characters or unitarity of \(\sigma\). These omissions prevent a fully convention-independent formulation; no additional hypotheses are imposed here.
\subsection{Short English statement}
For every number field, split reductive group, cuspidal automorphic representation, and finite-dimensional dual-group representation in the supplied scope, the conjecture asks for meromorphic continuation of the associated automorphic \(L\)-function to \(\C\) and the stated functional equation. Local factors, completion, and normalization conventions require clarification. The supplied status is partially solved; this is a statement-only draft requiring human mathematical review.
\subsection{Sources}
\begin{itemize}
\item Local UnsolvedMath registry: \texttt{data/unsolvedmath/problems.json}, permanent ID \(30004782\), source code \texttt{catalogue entry}, title ``Analytic Continuation and Functional Equations of Automorphic L-Functions,'' in the Oberwolfach Reports Open Problems collection. The supplied record lists proposed year 2021 and names no proposer.
\item Supplied report citation: \emph{Automorphic Forms, Geometry and Arithmetic (2022)}, Oberwolfach Reports, DOI 10.4171/owr/2021/39, pp.\ 2089--2156. \url{https://doi.org/10.4171/owr/2021/39}. The literature-source list labels this reference ``Oberwolfach Report 39/2021.'' The supplied background describes these reports as documenting mathematical workshops and their open-problem discussions. The differing year labels are retained as supplied.
\item Supplied literature reference: Braverman--Kazhdan, \emph{Gamma-functions of representations and lifting}. \url{https://arxiv.org/abs/math/9912208}.
\item Supplied literature assessment dated exactly \texttt{2026-08-22}: status ``partially solved.'' The supplied assessment states that meromorphic continuation and functional equations are established for \(\operatorname{GL}_n\) and many further families by Langlands--Shahidi, Rankin--Selberg, and functorial methods, while the extracted general automorphic \(L\)-function conjecture remains open. This is supplied dated curation derived from the OWR material, not an independently verified literature assessment; its time-sensitive status and cited claims require independent verification.
\item Supplied verification metadata: the statement is labeled \texttt{corrected\_verified}, with a claim that the standalone statement and hypotheses were checked against the official Oberwolfach report. That verification claim is reproduced as source metadata and has not been independently checked for this draft.
\end{itemize}
\end{document}
